% !TeX root = ../papers/nc-p3.tex
\section{左右の積の照合による混標数環帯の全次数認証}
\label{sec:rank-fifteen-overlap-annulus}
\begin{theorem}\label{thm:rank-fifteen-overlap-annulus}
標数零の通常な階数15窓に剰余標数3の付値を取り、
$0<\nu(y)<3\nu(x)$ かつ $5\nu(y)\le6\nu(3)$ とする。
この窓は標準Koszul左右AS正則で、各始点で四中心四次式の多項式環上
自由階数256である。
\end{theorem}
\begin{proof}
第\ref{sec:rank-fifteen-variable-zero-disk}節の表示で $\bar z=1+i$、
$\nu(H)=\nu(y)/6$ と取る。$R=\mathbb Z[z]_{(3,z^2+z+2)}[[H]]$ とし、普遍表示の分数体上の商の中で、
各長さと文字について語が生成する $R$ 格子を取る。
定理\ref{thm:rank-fifteen-sixth-center-minors}までの計算では、未切断経路の長さ5まで、
および四中心商の長さ6まで、必要な $H$ の冪で割った後に全主元が単元になる。
よって各格子を目標次元個の語で生成でき、分数体上の下界と合わせて自由基底となる。
$H$ 除算後の単元判定は有限精度で済むため、同じ消去を形式級数で持ち上げられる。
各格子は任意の小さい $H\ne0$ への特殊化後も、実際の語空間に全射する。

まず $R$ を3で還元し、十五桁を使って計算する。
四次中心の自由座標は標数零と3の双方で最後の始点の同じ語基底の同じ列で固定する。
未切断経路の剰余消去から得た生成列は、対応する還元格子を生成する。
実際、関係を課した後も、還元格子が自由で $H$ 捩れを持たないため、
計算に使う $H$ 除算後の行もその格子で関係を与える。
列数と階数が等しければこれらの列は基底であり、左右の積はその基底での還元に一致する。
中心方程式は単元主元で解け、前記の自由座標が一致するので、選んだ四式も一致する。
同じ議論は四式を課した格子に対して長さ6まで成り立つ。

四式を課した商を $B$ とする。まず $n\le5$ では、未切断経路の基底で
中心式と長さ $n-4$ の末尾との積を消去し、両側からの生成元の積の表を作る。
長さ5で四桁を使い、七桁残る。次に始点 $p$、文字 $h$ ごとに
末尾別の $B_{p,n-1}$ の四つの生成系と、先頭別の $B_{p+1,n-1}$ の
四つの生成系を直和に並べる。長さ $n-2$ の中間の生成元 $w$ と
先頭末尾の文字 $a,b$ に対して、結合則
\[
                       (aw)b-a(wb)=0
\]
の行を入れる。両項の座標は前段の左右の積表から得られる。
この直和は実際の商へ全射し、各行はその核に入る。消去後の自由列を
生成系に取り、元の直和の各列の像を次の両積表に用いる。
各 $H$ 除算は有効桁を一つ消費する。従って途中で余分な自由列や未使用の関係があっても
得られる次元は上界となる。

\path{computations/verify_rank_fifteen_overlap_annulus.py} は各 $H$ 係数を
$\mathbb F_3(z)$ の有理関数として計算する。除算前の定数係数はこの体で
厳密に零であり、単元主元の分子・分母は $z^2+z+2$ で還元して非零である。
長さ6から13で各段の最大精度損失は全始点で
\[
                             2,2,0,2,0,0,0,0
\]
となり、長さ6の後に五桁、最後に一桁が残る。ここで仮定は $3\in H^5\mathcal O$ と同値である。
長さ6までの五桁の表は上の自由格子の積表の還元なので、
$R\to\mathcal O/(H^5)$ により所定の窓の積表に移せる。
以後の除算はこの五桁だけで完了する。座標が $H$ を法として零なら、
実際の関係を分数体で $H$ で割って一桁少ない合同を得る。各始点における商の次元上界の列は
$(1+t+t^2+t^3)^4$ の係数列で、次数13で零になる。
四中心式の抽象的多項式環上の自由加群からの全射と一般点からの下界が
全次数で一致するので自由性が従う。定理\ref{thm:rank-fifteen-cubic-zero-disk}と
同じ左右Koszul複体の議論で正則性を得る。
\end{proof}

\begin{corollary}
数体 $L$ 内の $x,y$ が通常な窓を与え、$D=x^3/y$ が
3の上の全ての素点で正付値を持つとする。
\[
  5\operatorname{ord}_3\!\left(N_{L/\mathbb Q}(y)\right)\le6[L:\mathbb Q]
\]
ならば、その窓にも定理の全次数の結論が成り立つ。
\end{corollary}
\begin{proof}
各素点 $w\mid3$ で $\nu_w(3)=1$ と正規化する。
$\operatorname{ord}_3(N(y))/[L:\mathbb Q]$ は $\nu_w(y)$ の
重み $[L_w:\mathbb Q_3]/[L:\mathbb Q]$ の平均である。
よって $\nu_w(y)\le6/5$ となる $w$ がある。
この値が正なら $\nu_w(D)>0$ より本定理が適用できる。
零なら $\nu_w(x)>0$ なので定理\ref{thm:rank-fifteen-residue-line}、
負なら $\nu_w(y)<2\nu_w(x)$ なので定理\ref{thm:rank-fifteen-pole-cone}を使う。
\end{proof}
